% !TEX TS-program = lualatex
Конечной целью алгоритма TD-Gammon является получение численной оценки текущего положения на доске. Оценив все возможные состояния, достижимые из данной позиции, можно выбрать наилучшее --- то, которое обеспечивает наивысшую ожидаемую выгоду. Таким образом, задача выбора следующего хода сводится к регрессионной проблеме, в рамках которой модель предсказывает числовую оценку состояния на основе признаков текущей игровой ситуации.

Для решения этой задачи наиболее эффективной оказывается модель многослойного перцептрона (MLP) \cite{mlp}, поскольку она способна моделировать сложные нелинейные зависимости между конфигурацией на доске и соответствующими действиями.

Выход модели <<многослойный перцептрон>> вычисляется с помощью прямого распространения активаций: от входного слоя $\bar{l_1}$ к выходному слою $\bar{l_k}$, проходя через один или несколько ($k - 2$) скрытых слоёв $\bar{l_i}$. Каждое соединение между нейронами параметризуется вещественным весом $w_{ij}$. Каждый узел сети вычисляет взвешенную линейную сумму поступающих на него входов, после чего применяется нелинейная функция активации --- обычно сигмоидная функция (\ref{sigmoid}), которая сжимает результат в диапазон от 0 до 1 (\ref{sum-layer}).

Схематическое изображение структуры многослойного перцептрона представлено на рисунке \refimage{mlp-schema}

\image{Структура многослойного перцептрона}{mlp-schema}{mlp-schema}

\begin{equation}\label{sigmoid}
    \sigma(x) = \frac{1}{1 + e^{-x}}.
\end{equation}

\begin{equation}\label{sum-layer}
    \bar{l_j}= \sigma \left(\sum_i w_{ij} \bar{l_{j-1}} \right).
\end{equation}

где $\bar{l_j}$ --- вектор активаций на $j$-м слое сети, т.е. выходы нейронов этого слоя; $\bar{l_{j-1}}$ --- вектор активаций на предыдущем $(j-1)$-м слое; $w_{ij}$ --- вес, который соединяет $i$-й нейрон предыдущего слоя с $j$-м нейроном текущего слоя.
